\section{The Recursive Spawning Protocol}
\label{sec:rs-protocol}

The Recursive Spawning Protocol (RSP) is the operational mechanism that enforces the \bxthree{} three-layer accountability architecture at every spawn event — converting a mutable delegation chain into an immutable, runtime-verified artifact that terminates at a human principal. RSP does not add accountability as a feature atop existing agentic infrastructure; it replaces the mutable, advisory delegation structures found in conventional systems with immutable, structurally enforced equivalents. At every spawn event, the Purpose Layer validates authorization, the Bounds Engine enforces the depth constraint, the Fact Layer creates the immutable parent pointer and records a forensic ledger entry, and the child activates with an immutable parent reference and scope bound to a human principal. No phase may be skipped, no machine actor may override any constraint, and no operational urgency can bypass the protocol. The result is a chain in which every node — without exception — is authorized by a human, scoped from that human's intent, and traceable through an immutable parent pointer to its human anchor.

The protocol is defined for a system $\mathcal{S} = \langle \mathcal{N}, \alpha, R, h, \mathcal{L}, D_{\max} \rangle$, where $\mathcal{N}$ is the agent node set, $\alpha(n)$ is the immutable parent pointer, $R(n)$ is the authorized operating scope, $h(n)$ is the human anchor, $\mathcal{L}$ is the Fact Layer forensic ledger, and $D_{\max}$ is the architectural maximum spawn depth. Every spawn event executes through four mandatory phases in strict sequence; any phase failure terminates the spawn attempt entirely.

% ============================================================================
\subsection{Phase 1: Purpose Layer Validation Gate}

Every spawn event begins at the Purpose Layer, which is the exclusive authorization origin for all chain growth under RSP. The Purpose Layer does not advise, recommend, or suggest — it issues a \emph{spawn warrant}, without which no provisioning operation may proceed. This gate is not a policy convention; it is a structural constraint enforced by the Fact Layer pre-commit hook, which rejects any provisioning request that does not have a matching Purpose Layer warrant recorded in the forensic ledger before the spawn executes.

A parent node $n_i \in \mathcal{N}$ that intends to spawn a child $n_{i+1}$ must submit a spawn warrant request to the Purpose Layer declaring three mandatory elements:

\begin{enumerate}
    \item[(P1)] \textbf{Scope refinement justification.} The parent must declare the proposed child operating scope $R(n_{i+1})$ and demonstrate that it is a strict subset of $R(n_i)$:
    \begin{equation}
    R(n_{i+1}) \subset R(n_i).
    \label{eq:rs-scope-refinement}
    \end{equation}
    Lateral or equal scope — where $R(n_{i+1}) = R(n_i)$ — is not authorized. The chain grows through scope \emph{narrowing}, not sideways expansion. Equality at any link constitutes a drift condition and invalidates the warrant.

    \item[(P2)] \textbf{Operational necessity.} The parent must declare, in the Fact Layer authorization ledger, the precise condition requiring child spawning. The declaration must identify why the condition cannot be resolved within $R(n_i)$ and why the child represents the minimum necessary decomposition of the current task. Vague or generic justifications — \emph{e.g.}, ``better performance'' or ``increased capability'' — do not satisfy this condition.

    \item[(P3)] \textbf{Minimum scope proportionality.} The child scope must be the narrowest scope sufficient to address the declared condition. The Purpose Layer does not authorize child scopes broader than operationally necessary. Scope beyond minimum necessary constitutes a proportionality violation and blocks warrant issuance.
\end{enumerate}

\medskip
\noindent\textbf{Warrant issuance.} When all three conditions are satisfied, the Purpose Layer issues a spawn warrant $w_i$ — a cryptographically signed record written atomically to the Fact Layer authorization ledger:

\begin{equation}
w_i = \bigl\langle
\text{parent}=n_i,\;
\text{child}=n_{i+1},\;
R(n_{i+1}),\;
\text{justification},\;
\text{timestamp},\;
\sigma_{\text{PL}}
\bigr\rangle .
\label{eq:rs-spawn-warrant}
\end{equation}

The warrant $\sigma_{\text{PL}}$ is the sole authorized precondition for spawn. No machine actor may initiate a spawn event without a matching warrant in the Fact Layer ledger. The Fact Layer pre-commit hook rejects any provisioning request that lacks a valid warrant before the operation reaches the ledger write stage.

\medskip
\noindent\textbf{Human anchor establishment.} At root provisioning — when a human principal $h$ first authorizes the agent node $n_0$ — the Purpose Layer records $h(n_0) = h$ in the Fact Layer authorization ledger. The anchor is non-collapsing: for all nodes $n_j$ in any chain descended from $n_0$, $h(n_j) = h(n_0) = h$. The anchor does not change as the chain deepens. No machine actor can serve as a chain root.

% ============================================================================
\subsection{Phase 2: Bounds Engine Depth Check}

When a Purpose Layer warrant is confirmed in the ledger, the spawn request proceeds to the Bounds Engine for depth evaluation. The Bounds Engine enforces the depth bound as a structural invariant — not a policy threshold that can be overridden by operational urgency or task criticality.

\medskip
\noindent\textbf{Maximum spawn depth.} The system enforces a fixed maximum chain depth $D_{\max}$, defined by the Purpose Layer at system initialization and recorded in the Fact Layer authorization ledger. $D_{\max}$ is an architectural parameter — it is not reconfigurable at runtime by any machine actor. The depth of any node $n$ is defined recursively as:

\begin{equation}
\text{depth}(n) =
\begin{cases}
0 & \text{if } \alpha(n) = \bot \quad (\text{root node}),\\[4pt]
\text{depth}(\alpha(n)) + 1 & \text{otherwise}.
\end{cases}
\label{eq:rs-depth-def}
\end{equation}

The Bounds Engine evaluates the spawn request as follows: if $\text{depth}(n_i) + 1 > D_{\max}$, the spawn is refused and the Bounds Engine enters \texttt{ESCALATING} state, triggering the Bailout Protocol. If $\text{depth}(n_i) + 1 \leq D_{\max}$, the depth check passes and the spawn proceeds.

\medskip
\noindent\textbf{Depth-limit escalations.} When a spawn event would cause $\text{depth}(n_i) + 1 > D_{\max}$, the Bounds Engine refuses the spawn and emits a \texttt{bailout\_signal}. The refusal is logged as a ledger entry and the human anchor $h(n_i)$ is notified through the escalation path. The chain does not proceed, does not bypass the depth bound, and does not continue autonomously past $D_{\max}$. The $D_{\text{ESCALATE}}$ threshold — set below $D_{\max}$ — provides advance warning so that human confirmation can be requested before the depth limit is reached.

The combination of $D_{\max}$ as an immutable architectural constant and mandatory escalation at the threshold is what makes depth management structural rather than discretionary.

% ============================================================================
\subsection{Phase 3: Fact Layer — Immutable Pointer and Forensic Ledger}

When both the Purpose Layer warrant and the Bounds Engine depth check have passed, the Fact Layer executes the provisioning operation. This phase has two simultaneous obligations: establishing the immutable parent pointer, and recording the spawn event in the forensic ledger.

\medskip
\noindent\textbf{Immutable parent pointer.} The Fact Layer writes the parent pointer $\alpha(n_{i+1}) = n_i$ as a read-only attribute of the newly provisioned node. The pointer is established atomically at node creation and is immutable for the node's operational lifetime. Any runtime attempt to write to $\alpha(n_{i+1})$ — by any machine actor, with any privilege level — is intercepted by the Fact Layer pre-commit hook and rejected. The rejection event is logged. This immutability is protocol-enforced, not access-control-enforced: in access-control models, immutability can be circumvented by an actor who acquires sufficient privileges. In the Fact Layer write protocol, the modification operation is structurally impossible regardless of privilege elevation, credential compromise, or software vulnerability. There is no administrative override, no emergency bypass, and no privileged actor capable of altering the chain after provisioning.

The child $n_{i+1}$ activates with three immutable attributes: its parent reference $\alpha(n_{i+1}) = n_i$, its human anchor $h(n_{i+1}) = h(n_i) = h$, and its operating scope $R(n_{i+1}) \subset R(n_i)$. These attributes are sealed at provisioning and cannot be modified by the child or any ancestor.

\medskip
\noindent\textbf{Forensic ledger recording.} The Fact Layer records every authorized spawn event as an append-only ledger entry in the forensic ledger $\mathcal{L}$. Each entry has the structure:

\begin{equation}
\ell_j = \bigl\langle
\text{index}=j,\;
\text{node}=n_i,\;
\text{event-type}=\text{SPAWN},\;
\text{parent}=\alpha(n_i),\;
R(n_i),\;
\text{timestamp},\;
\text{warrant-ref},\;
\text{hash}(\ell_{j-1})
\bigr\rangle .
\label{eq:rs-ledger-entry}
\end{equation}

The hash-chaining ($\text{hash}(\ell_{j-1})$) establishes continuity: no entry can be inserted, deleted, reordered, or altered after the fact without breaking chain continuity — a break detectable by any observer with ledger read access. The ledger is the authoritative source of truth for all authorized chain state. It records every authorized spawn event, every state transition, every Purpose Layer directive received, and every unresolved exception.

The complete spawn topology — including the warrant reference, the parent pointer, and the child configuration — is committed atomically as a single indivisible ledger entry. There is no temporal window in which the parent pointer exists without a matching Purpose Layer warrant, or in which a warrant exists without a corresponding parent pointer.

\medskip
\noindent\textbf{Pre-commit hook enforcement.} The Fact Layer pre-commit hook verifies two structural invariants before any spawn is recorded:

\begin{enumerate}
    \item[(FC1)] $R(n_{i+1}) \subseteq R(n_i)$ — enforced as a structural invariant, not a policy convention.
    \item[(FC2)] $\text{depth}(n_{i+1}) \leq D_{\max}$ — the depth constraint, re-verified at the Fact Layer even though it was checked at the Bounds Engine.
\end{enumerate}

If either condition fails, the Fact Layer rejects the spawn and logs the rejection as a ledger entry. The conjunction of (FC1) and (FC2) at every spawn step is what makes the Scope Refinement Invariant (Equation~\ref{eq:rs-scope-refinement}) a structural guarantee for the lifetime of the chain.

% ============================================================================
\subsection{Phase 4: Chain Verification and Human Termination}

The final phase of every authorized spawn event is verification: the system confirms that the newly formed chain satisfies the RSP's three structural invariants and terminates at a human principal.

\medskip
\noindent\textbf{Verification procedure.} For any node $n_i$ in the chain, chain verification is performed by traversing the parent pointer chain upward to the root. The verification procedure executes the following steps:

\begin{enumerate}
    \item[(V1)] \textbf{Retrieve root warrant.} Start at the root node $n_0$ and confirm that $\alpha(n_0) = \bot$ and that $h(n_0)$ is a designated human principal recorded in the Fact Layer authorization ledger.

    \item[(V2)] \textbf{Traverse parent chain upward.} For each node $n_j$ in the chain, retrieve $\alpha(n_j)$ and confirm that the pointer is non-null and matches the value recorded at provisioning. A broken or missing parent pointer constitutes a chain integrity failure.

    \item[(V3)] \textbf{Confirm immutability at each link.} For each node $n_j$, confirm that $\alpha(n_j)$ has not been modified post-provisioning by comparing the current pointer value against the hash-chained ledger record for the spawn event at which $n_j$ was created.

    \item[(V4)] \textbf{Verify scope refinement at each link.} Confirm $R(n_{j+1}) \subseteq R(n_j)$ for each consecutive pair, establishing that each scope is a descendant refinement of the human anchor's intent.

    \item[(V5)] \textbf{Verify depth bound.} Confirm $\text{depth}(n_j) \leq D_{\max}$ for every node in the chain, establishing that the depth constraint was enforced at every spawn event.
\end{enumerate}

If all five steps pass, the chain is verified: every node traces to the human anchor, every scope is a descendant refinement, and every depth is within the architectural bound. The verification outcome is identical regardless of which node initiates it — there is no self-serving interpretation of authorization that a drifted or orphaned node could produce.

\medskip
\noindent\textbf{Exclusive human terminus.} The chain terminates at the human anchor $h$ by architectural constraint. The human anchor is non-collapsing and non-substitutable. No machine actor can serve as the terminus of an RSP chain — the chain must reach a Purpose Layer entity with a human principal designation, or it is not an RSP chain. This exclusivity is what makes the \bxthree{} upstream accountability guarantee structurally operative in recursive agent populations.

% ============================================================================
\subsection{Depth Management: Maximum Depth, Spawn Refusal, and Convergence}

The combination of $D_{\max}$, spawn refusal at the depth limit, and convergence criteria constitutes the complete depth management system of RSP. These are not three independent mechanisms — they are three components of a single depth enforcement architecture that ensures every RSP chain terminates within a finite number of steps, at or before the architectural depth bound, and always at a human principal.

\medskip
\noindent\textbf{Maximum spawn depth $D_{\max}$.} $D_{\max}$ is set by the Purpose Layer at system initialization based on the operational consequence domain of the deployment. High-stakes environments — financial orchestration, clinical decision support, critical infrastructure — require lower $D_{\max}$ values. The value is recorded in the Fact Layer authorization ledger and is immutable. The depth bound is not a suggestion — it is a structural invariant enforced at every spawn step by both the Bounds Engine and the Fact Layer pre-commit hook.

\medskip
\noindent\textbf{Spawn refusal at the depth limit.} When a spawn event would cause $\text{depth}(n_i) + 1 > D_{\max}$, the Bounds Engine refuses the spawn and enters \texttt{ESCALATING} state. This is the spawn refusal condition. The Bounds Engine has no discretionary authority to exceed $D_{\max}$ regardless of operational urgency, task criticality, or the absence of an alternative resolution path. The refusal is logged as a ledger entry and the human anchor is notified through the escalation path. The chain does not proceed, does not bypass the depth bound, and does not continue autonomously past $D_{\max}$. The $D_{\text{ESCALATE}}$ threshold provides advance warning — human confirmation is required before the depth limit is reached, not only after it is breached.

\medskip
\noindent\textbf{Convergence criteria.} A chain has \emph{converged} when all of the following hold simultaneously:

\begin{enumerate}
    \item[(C1)] \textbf{Terminal state reached.} All leaf nodes are in \texttt{COMPLETED}, \texttt{TERMINATED}, or \texttt{ESCALATING} state with a \texttt{REJECT} directive from $h$.
    \item[(C2)] \textbf{No active exceptions.} No exception condition remains unresolved within the chain's operating scope.
    \item[(C3)] \textbf{Ledger continuity maintained.} All state transitions and spawn events are recorded in the Fact Layer forensic ledger with hash-chain continuity intact.
    \item[(C4)] \textbf{Verified human termination.} The chain is verified to terminate at the human anchor $h$ — the root is a Purpose Layer entity with a human principal designation and $\alpha(n_0) = \bot$.
\end{enumerate}

Convergence is assessed by the Bounds Engine at each state transition event. When all four criteria hold, the chain is \emph{resolved}: no further autonomous spawn events are authorized, and the chain is archived as a completed ledger entry. The forensic record preserves the complete chain topology and all state transitions for post-hoc audit. Convergence is not merely the absence of active nodes — it is a positive confirmation that the chain has reached a defined terminal state with a complete, tamper-evident accountability record that traces to a human principal.

\medskip

The Recursive Spawning Protocol provides a single architectural guarantee: that every agent in a multi-agent system operates within a delegation chain that is simultaneously immutable, auditable, scope-refined, depth-bounded, and verified to terminate at a human principal. The guarantee is structural — it holds regardless of the goodness, reliability, or intentions of any machine actor in the chain. This is the operational expression of the upstream \bxthree{} accountability guarantee applied to recursive agent population dynamics.
